\section{Component Design}

\subsection{Sensors}

Sensors provide temperature samples in millidegrees Celsius. Platform backends may read Linux \texttt{hwmon}, files, serial HIL input, 1-Wire probes, I2C sensors, or future automotive sensor paths. The core receives only numeric sensor IDs and typed values.

\subsection{Actuators}

The v1 actuator model is a PWM fan with optional tach feedback. The actuator definition includes minimum PWM, maximum PWM, slew limit, and stall threshold.

\subsection{Context Inputs}

Vehicle speed is the first context input. It is deliberately not special in the core: future context inputs such as ignition state, ambient temperature, HVAC state, and drive mode should use the same mechanism.

\subsection{Fault Detectors}

Initial fault detectors:

\begin{itemize}[leftmargin=*]
    \item fan stall while PWM is above a configured threshold;
    \item stuck sensor over a configured observation window;
    \item thermal runaway when temperature rises despite requested cooling;
    \item stale context input --- for example a speed-source timeout, which drives acoustic-mask fail-safe behavior.
\end{itemize}

\subsection{Numeric Identity}

Every entity the core manipulates is addressed by a small integer, never a string: sensors, actuators, zones, and context signals carry numeric IDs assigned at configuration load. Telemetry signal IDs are partitioned into fixed namespaces --- zone temperatures, actuator outputs, PID terms, context values, modifier outputs, and fault counters --- so a decoder can classify a signal from its ID alone. Fault detectors instantiate per target: one stall detector per actuator, one stuck-sensor detector per sensor, one runaway detector per zone, one stale-context detector per context signal. Each detector runs a common state machine --- normal, then degraded or critical or latched, then recovering --- with per-detector severity, action, and persistence and recovery tick counts. The exact struct layouts, enum values, and the signal-ID and event-code tables are project header artefacts and are catalogued in the appendix reference rather than transcribed here, so the paper cannot drift out of sync with the code.

\endinput
